\section[Appendix - Use of Generative AI in Research]{Use of Generative AI in Research}\label{app:genai}

This appendix documents what Generative AI (GenAI) tools were used in this
research, how they were deployed, and what functions they served.

Four GenAI systems supported this research:

\begin{enumerate}
    \item Perplexity: For literature discovery and initial source identification
    \item Anthropic Claude: For critical analysis of reading notes, validation of
          guidance artifact outputs, generation of the machine-readable YAML
          companion,\footnote{\url{https://github.com/thiagoavadore/thesis_2024-2026_antwerp_thesis_GenAIPracticeVacuum/blob/main/guidance-artifact.yaml}} front cover prompt design,
          and thesis revision assistance
    \item Meta Llama 3 (self-hosted instance): For critical analysis of reading notes
    \item Google Gemini: For front cover image generation
\end{enumerate}

These tools served six functions:

\begin{enumerate}
    \item \textbf{Efficiency in discovery}: Accelerated identification of relevant
          sources across disciplines;
    \item \textbf{Critical examination}: Provided external perspective on blind spots
          and gaps in interpretation;
    \item \textbf{Research alignment}: Helped assess literature relevance to evolving
          research questions;
    \item \textbf{Artifact validation}: Checked practitioner checklist cards for
          internal consistency, completeness against Section~\ref{sec:artifact}
          content, and formatting;
    \item \textbf{Format generation}: Produced the machine-readable YAML guidance
          artifact from researcher-authored content;
    \item \textbf{Image generation}: Produced the front cover image from a
          researcher-designed concept.
\end{enumerate}

The following sections document each usage in detail.

\subsection{Literature Discovery Process}

Perplexity helped identify relevant literature beyond traditional academic
search engines (Google Scholar) and supervisor recommendations.

\textbf{Example discovery prompt}: \textit{In the context of Software and Infrastructure engineering,
    are there any peer-reviewed articles/papers that connect Reference Architecture and tools that
    support it, including AI tools? If so, give me a list with sources. This is for academic research,
    so explore and reference only journals, thesis and dissertations - skip the blog posts.}

\subsection{Literature Review Process}

The review process followed a three-stage pattern:

\begin{enumerate}
    \item \textbf{Reading and Note-Taking}: Papers identified through discovery were added to
          Zotero and read systematically. Notes captured key points, learnings, questions,
          and findings.
    \item \textbf{Critical Analysis}: Notes were submitted to GenAI systems for critical
          examination using the following prompt:

          \textit{I've read the paper (attached) and took notes (also attached) by trying to summarize
              interesting points, learnings, questions and findings. Investigate potential blind spots,
              contradictions or any type of hand-waving on my notes. Critique my points and avoid
              agreeing with me for the sake of making me happy. If the article publication date is before
              your cut training date, investigate what theories and implications it had that are not
              highlighted in my notes, always giving me link to sources.}

    \item \textbf{Research Alignment}: Once the Problem Statement and Research Questions
          crystallized, the analysis prompt evolved to include research alignment:

          \textit{As part of my Executive MSc I'm researching into this Problem Statement/Research
              Question (attached). To better investigate it, I've read the paper (attached) and I
              took notes (also attached) by trying to summarize interesting points, learnings, questions
              and findings. Investigate potential blind spots, contradictions or any type of hand-waving
              on my notes. Critique my points and avoid agreeing with me for the sake of making me happy.
              If the article publication date is before your cut training date, investigate what theories
              and implications it had that are not highlighted in my notes, always giving me link to sources.
              And lastly, investigate if the article supports, contradicts or adds new information to
              the Problem Statement or the Research Question.}

\end{enumerate}

\subsection{Guidance Artifact Preparation}

Two outputs in the appendices involved GenAI assistance:

\begin{enumerate}
    \item \textbf{Practitioner Checklist Cards (\ref{app:checklists}):} The checklist cards were
          researcher-authored based on Delphi findings and the guidance artifact in
          Section~\ref{sec:artifact}. GenAI (Claude) was used to validate the cards for internal
          consistency, completeness against Section~\ref{sec:artifact} content, and formatting. GenAI
          did not generate the card content.
    \item \textbf{YAML Guidance Artifact:} The machine-readable YAML companion
          (available in the research repository\footnote{\url{https://github.com/thiagoavadore/thesis_2024-2026_antwerp_thesis_GenAIPracticeVacuum/blob/main/guidance-artifact.yaml}}) was generated by GenAI (Claude) from the
          researcher-authored Section~\ref{sec:artifact} and checklist card content. The researcher
          reviewed, corrected, and approved the output. This constitutes a format
          transformation of existing researcher-authored material, not new content
          generation.
\end{enumerate}

\subsection{Thesis Revision}

During the final revision phase, Claude (Anthropic) assisted with mechanical
verification tasks: identifying misformatted cross-references after structural
reorganization and verifying APA citation consistency. In all cases, the
researcher reviewed the identified issues and applied corrections manually.

\subsection{Front Cover Image}

The researcher designed the front cover concept to represent the thesis theme
visually: cognitive work (represented by an anatomical brain) transforming into
an architectural blueprint through a verification lens. Three design constraints
guided the concept. First, the image had to be monochromatic to avoid dominating
the page. Second, the transformation had to show intentional loss and change in
the process, reflecting how GenAI outputs require curation rather than blind
acceptance. Third, the style had to match the scholarly tone of an academic
thesis.

Anthropic Claude translated these constraints into a structured image generation
prompt. Google Gemini then generated the image from the following prompt:

\textit{Monochromatic grayscale academic illustration, no text, no letters, no
words, no numbers. A detailed anatomical human brain on the left side, drawn in
fine ink lines, partially dissolving into geometric particles and dots that drift
rightward. These particles gradually reassemble into an abstract architectural
blueprint floorplan on the right side, shown as precise technical line drawings
of rooms, corridors, and structural elements. Between the brain and the
blueprint, a translucent magnifying lens floats, refracting the particles
passing through it, symbolizing verification and curation. Subtle watercolor ink
splashes in light gray behind both elements. Small geometric shapes (circles,
hexagons, connection nodes) scattered in the transition zone. Style: scientific
illustration meets architectural drawing, similar to academic thesis cover art.
White background, no color, purely grayscale ink and pencil aesthetic. High
detail, scholarly tone.}

\subsection{What GenAI Did Not Do}

GenAI tools did not:

\begin{enumerate}
    \item Generate original thesis argumentation or theoretical analysis;
    \item Make methodological decisions;
    \item Conduct analysis of primary research data;
    \item Assist with between-round synthesis, including data coding, taxonomy
          development, or preparation of subsequent round materials;
    \item Develop the conceptual framework;
    \item Independently interpret supervisor feedback.
\end{enumerate}

All interpretation, synthesis, framework development, and argumentation is
the researcher's own work.
